\chapter{Reports \& Long Documents}
\label{ch:reports}

This chapter covers five \texttt{scrreprt}-based doctypes.  These document
types support chapters via the KOMA-Script report class, making them suitable
for multi-section technical works, proposals, and reference documents.  All
share \texttt{listof=totoc}, \texttt{bibliography=totoc},
\texttt{chapterprefix=false}, \texttt{open=any}, and
\texttt{numbers=noenddot} from the default \texttt{scrreprt} options.

\section{Common Features}

Every report-type doctype provides:

\begin{itemize}
    \item \textbf{Chapter-level sectioning} --- Full
          \texttt{\textbackslash chapter} command for top-level divisions,
          with sections and subsections below.
    \item \textbf{Open on any page} --- \texttt{open=any} allows chapters
          to start on even or odd pages, saving space.
    \item \textbf{Lists in TOC} --- \texttt{listof=totoc} automatically
          includes the list of figures, list of tables, and list of
          algorithms in the table of contents.
    \item \textbf{No chapter prefix} --- \texttt{chapterprefix=false}
          omits the ``Chapter\,N'' prefix from chapter headings.
\end{itemize}

\section{Choosing the Right Report Type}

\begin{table}[htbp]
    \centering
    \caption{Report-type doctype decision guide}
    \label{tab:report-decision}
    \begin{tabular}{@{} l p{6.5cm} @{}}
        \toprule
        \textbf{Doctype} & \textbf{Best For} \\
        \midrule
        report / technical-report & Internal or client-facing technical documents with report numbers \\
        research-proposal         & Grant applications with PI, budget, and programme metadata \\
        standard                  & Formal standards documents with designation and ICS codes \\
        patent                    & Patent specifications with legal-style sectioning \\
        manual                    & Product manuals, guides, or handbooks \\
        \bottomrule
    \end{tabular}
\end{table}

\section{Doctype Profiles}

\subsection{report / technical-report}

The \texttt{technicalreport} doctype (aliased as \texttt{report},
\texttt{tech-report}, etc.) produces a formal technical report with a title
page containing report number, revision, author, client, project name,
confidentiality level, and a personnel table.  It is single-sided with
chapter-level sectioning.

\begin{table}[htbp]
    \centering
    \caption{Custom commands for \texttt{technicalreport}}
    \label{tab:techreport-commands}
    \begin{tabular}{@{} l p{8cm} @{}}
        \toprule
        \textbf{Command} & \textbf{Description} \\
        \midrule
        \texttt{\textbackslash reportnumber\{...\}}    & Report identifier (default: TR-000) \\
        \texttt{\textbackslash revision\{...\}}        & Revision letter (default: A) \\
        \texttt{\textbackslash preparedby\{...\}}      & Preparing group or author \\
        \texttt{\textbackslash client\{...\}}          & Client organisation \\
        \texttt{\textbackslash projectname\{...\}}     & Associated project name \\
        \texttt{\textbackslash confidentiality\{...\}}  & Classification (default: Internal Use Only) \\
        \texttt{\textbackslash sponsor\{...\}}         & Executive sponsor name \\
        \texttt{\textbackslash leadengineer\{...\}}    & Lead engineer name \\
        \texttt{\textbackslash doccontrol\{...\}}      & Document control contact \\
        \texttt{\textbackslash distribution\{...\}}    & Distribution statement \\
        \bottomrule
    \end{tabular}
\end{table}

Use \texttt{technicalreport} for engineering reports, consultant
deliverables, or any document that tracks report numbers and revisions.

\begin{minted}{latex}
\documentclass[doctype=technical-report]{../../omnilatex}
\title{Structural Analysis of Bridge Deck Section B-7}
\reportnumber{TR-2025-042}\revision{C}
\preparedby{Structural Analysis Group}
\client{Department of Transportation}
\projectname{Highway 101 Overpass Rehabilitation}
\confidentiality{Public Release}
\begin{document}
\maketitle
\tableofcontents
\chapter{Executive Summary}
This report presents the findings of the structural analysis.
\end{document}
\end{minted}

\subsection{research-proposal}

The \texttt{research-proposal} doctype is tailored for grant and fellowship
applications.  The title page displays the funding programme, call reference,
principal and co-investigators, organisation, duration, budget, and start/end
dates.  It uses \texttt{chapterprefix=false} and \texttt{open=any} for a
compact layout.

\begin{table}[htbp]
    \centering
    \caption{Custom commands for \texttt{research-proposal}}
    \label{tab:proposal-commands}
    \begin{tabular}{@{} l p{8cm} @{}}
        \toprule
        \textbf{Command} & \textbf{Description} \\
        \midrule
        \texttt{\textbackslash propabstract\{...\}}         & Proposal abstract or summary \\
        \texttt{\textbackslash propprogram\{...\}}          & Funding programme name \\
        \texttt{\textbackslash propcall\{...\}}             & Call for proposals identifier \\
        \texttt{\textbackslash propinvestigator\{...\}}     & Principal Investigator name \\
        \texttt{\textbackslash propcoinvestigators\{...\}}  & Co-Investigator names \\
        \texttt{\textbackslash proporganization\{...\}}     & Host institution or organisation \\
        \texttt{\textbackslash propduration\{...\}}         & Project duration (default: 24 months) \\
        \texttt{\textbackslash propbudget\{...\}}           & Total requested budget \\
        \texttt{\textbackslash propstartdate\{...\}}        & Project start date \\
        \texttt{\textbackslash propenddate\{...\}}          & Project end date \\
        \bottomrule
    \end{tabular}
\end{table}

Use \texttt{research-proposal} for ERC, NSF, DFG, or similar funding
applications that require structured metadata on the title page.

\begin{minted}{latex}
\documentclass[doctype=research-proposal]{../../omnilatex}
\title{Machine Learning for Predictive Maintenance}
\propprogram{Horizon Europe}
\propcall{HORIZON-CL4-2025-01}
\propinvestigator{Dr.\ Jane Doe}
\propcoinvestigators{Dr.\ A.~Smith, Dr.\ B.~Mueller}
\proporganization{Technical University}
\propduration{36 months}
\propbudget{EUR 1,500,000}
\propstartdate{2026-01-01}\propenddate{2028-12-31}
\begin{document}
\maketitle
\chapter{Project Description}
This proposal aims to develop ML-based predictive maintenance.
\end{document}
\end{minted}

\subsection{standard}

The \texttt{standard} doctype is designed for formal standards documents
resembling ISO or IEC publications.  The title page shows the standards
series, designation number, status, scope, committee, publication date,
superseded document, and ICS codes.  Section numbering uses plain Arabic
numerals without chapter prefixes (e.g.\ 1, 1.1, 1.1.1).

\begin{table}[htbp]
    \centering
    \caption{Custom commands for \texttt{standard}}
    \label{tab:standard-commands}
    \begin{tabular}{@{} l p{8cm} @{}}
        \toprule
        \textbf{Command} & \textbf{Description} \\
        \midrule
        \texttt{\textbackslash standardseries\{...\}}         & Standards series name \\
        \texttt{\textbackslash standarddesignation\{...\}}    & Designation number (default: F-000) \\
        \texttt{\textbackslash standardstatus\{...\}}         & Status (Draft, Published, Withdrawn) \\
        \texttt{\textbackslash standardscope\{...\}}          & Scope statement displayed on title page \\
        \texttt{\textbackslash standardcommittee\{...\}}      & Preparing committee name \\
        \texttt{\textbackslash standardsupersedes\{...\}}     & Superseded standard designation \\
        \texttt{\textbackslash standardics\{...\}}            & ICS classification codes \\
        \texttt{\textbackslash standardkeywords\{...\}}       & Subject keywords \\
        \texttt{\textbackslash standarddisclaimer\{...\}}     & Legal disclaimer text \\
        \bottomrule
    \end{tabular}
\end{table}

Use \texttt{standard} for internal standards, specifications, or normative
documents that require designation tracking and formal status indicators.

\begin{minted}{latex}
\documentclass[doctype=standard]{../../omnilatex}
\title{Thermal Performance of Building Envelopes}
\standardseries{Omni Standards Series}
\standarddesignation{B-204}
\standardstatus{Published}
\standardscope{This standard specifies methods for measuring
  the thermal transmittance of building envelope components.}
\standardcommittee{TC 163 -- Building Environment}
\standardics{91.120.10}
\begin{document}
\maketitle
\section{Scope}
This standard applies to opaque and transparent building elements.
\end{document}
\end{minted}

\subsection{patent}

The \texttt{patent} doctype provides a minimal report-class layout for patent
specifications.  It sets \texttt{titlepage=true}, single-sided, with
sectioning depth limited to two levels (\texttt{secnumdepth=2},
\texttt{tocdepth=1}).  No custom commands are defined---the doctype relies
on standard KOMA-Script commands and the user's own preamble.

Use \texttt{patent} for drafting patent applications or publishing patent
descriptions where the formal structure (title page, sections, claims) is
composed directly by the author.

\begin{minted}{latex}
\documentclass[doctype=patent]{../../omnilatex}
\title{Method and Apparatus for Efficient Heat Recovery}
\author{Jane Doe}\date{May 2025}
\begin{document}
\maketitle
\chapter{Field of the Invention}
This invention relates to heat recovery systems for industrial processes.
\chapter{Background}
Existing heat recovery methods suffer from thermal losses exceeding 30\%.
\chapter{Summary}
The present invention provides a heat recovery apparatus with improved
thermal efficiency through a counter-flow heat exchanger design.
\end{document}
\end{minted}

\subsection{manual}

The \texttt{manual} doctype is optimised for product manuals, user guides,
and handbooks.  It uses a sans-serif font at 11\,pt with single line spacing
and \texttt{DIV=14} for comfortable reading.  The front page displays a brand
name, document title, firmware or software version, release date, support
contact, and a summary paragraph.  No separate title page is generated; the
front matter appears on the first page with \texttt{titlepage=false}.

\begin{table}[htbp]
    \centering
    \caption{Custom commands for \texttt{manual}}
    \label{tab:manual-commands}
    \begin{tabular}{@{} l p{8cm} @{}}
        \toprule
        \textbf{Command} & \textbf{Description} \\
        \midrule
        \texttt{\textbackslash manualbrand\{...\}}         & Product or brand name (default: OmniLaTeX Devices) \\
        \texttt{\textbackslash manualversion\{...\}}       & Firmware/software version (default: 1.0) \\
        \texttt{\textbackslash manualreleasedate\{...\}}   & Release date (default: \texttt{\textbackslash today}) \\
        \texttt{\textbackslash manualsupport\{...\}}       & Support contact information \\
        \texttt{\textbackslash manualsummary\{...\}}       & One-paragraph manual summary \\
        \bottomrule
    \end{tabular}
\end{table}

Use \texttt{manual} for hardware manuals, software documentation, or any
reference guide that needs version tracking and support contact information.

\begin{minted}{latex}
\documentclass[doctype=manual]{../../omnilatex}
\title{Model X-2000 User Guide}
\manualbrand{Acme Devices}
\manualversion{3.2.1}
\manualsupport{support@acme.com}
\manualsummary{This guide covers installation, configuration,
  and operation of the Model X-2000.}
\begin{document}
\maketitle
\tableofcontents
\chapter{Getting Started}
Unpack the device and verify all components against the parts list.
\end{document}
\end{minted}
